% ---- pg-344 （书页 332） ----

\begin{mvtable}
\phantomsection\label{tab:7.7}\mvtabcap{表 7.7\quad 纸浆与纸张特性数据}
\footnotesize
\begin{tabular}{cccccccc}
\toprule
$y_1$&$y_2$&$y_3$&$y_4$&$z_1$&$z_2$&$z_3$&$z_4$\\
BL&EM&SF&BS&AFL&LFF&FFF&ZST\\
\midrule
21.312&7.039&5.326&0.932&$-0.030$&35.239&36.991&1.057\\
21.206&6.979&5.237&0.871&0.015&35.713&36.851&1.064\\
20.709&6.779&5.060&0.742&0.025&39.220&30.586&1.053\\
19.542&6.601&4.479&0.513&0.039&39.756&21.072&1.050\\
20.449&6.795&4.912&0.577&$-0.070$&32.991&36.570&1.049\\
$\vdots$&$\vdots$&$\vdots$&$\vdots$&$\vdots$&$\vdots$&$\vdots$&$\vdots$\\
16.441&6.315&2.997&$-0.400$&$-0.605$&2.845&84.554&1.008\\
16.294&6.572&3.017&$-0.478$&$-0.694$&1.515&81.988&0.998\\
20.289&7.719&4.866&0.239&$-0.559$&2.054&8.786&1.081\\
17.163&7.086&3.396&$-0.236$&$-0.415$&3.018&5.855&1.033\\
20.289&7.437&4.859&0.470&$-0.324$&17.639&28.934&1.070\\
\bottomrule
\end{tabular}
\end{mvtable}

\noindent{\footnotesize 资料来源: 见文献[19].}

\medskip
\noindent (b) 用全部四个响应变量 $Y_1,Y_2,Y_3$ 和 $Y_4$ 以及四个自变量 $Z_1,Z_2,Z_3$ 和 $Z_4$ 作多元多重回归分析.

\noindent (i) 拟合并提出适当的回归模型. 给出估计的系数矩阵 $\hat{\boldsymbol{\beta}}$ 和估计的误差协方差矩阵 $\hat{\boldsymbol{\Sigma}}$.

\noindent (ii) 分析残差. 检查离群点.

\noindent (iii) 给定 (a) 中 (iii) 的自变量值, 对四个响应变量 $Y_1,Y_2,Y_3$ 和 $Y_4$ 构造 95\% 联合预测区间. 将此 $Y_3$ 的预测区间与 (a) 中 (iii) 所得到的预测区间作比较并解释.

\medskip
\noindent 7.27\quad 考虑\hyperref[tab:6.20]{表 6.20} 中关于手机继电塔的固定崩溃数据. 在最初的试验中, 经验水平记录为两个: 新手或专家. 现在考虑三个水平: 新手、专家和有经验者. 对某有经验的工程师又追加了一些试验, 数据如下. 另外, 在原来的数据中, 将运行 3 原来的分类专家改分为有经验者, 并将运行 14 原来的分类改分为有经验者. 这两组数据的其他数值保持不变. 以评估时间和解决问题时间为响应变量, 以问题严重性、问题复杂水平和工程师水平为预测变量, 作多元多重回归分析. 考虑包含预测变量和两因子交互效应的回归模型. (提示: 由于原来数据中的两个改变和新数据的加入使得设计变得不平衡, 因此建议用回归方法来分析.)

\begin{mvtable}
\begin{tabular}{lllccc}
\toprule
问题严重性&问题复杂水平&工程师水平&问题评估时间&问题解决时间&总修复时间\\
\midrule
低&复杂&有经验者&5.3&9.2&14.5\\
低&复杂&有经验者&5.0&10.9&15.9\\
高&简单&有经验者&4.0&8.6&12.6\\
高&简单&有经验者&4.5&8.7&13.2\\
高&复杂&有经验者&6.9&14.9&21.8\\
\bottomrule
\end{tabular}
\end{mvtable}

\phantomsection\label{bib:7}
\section*{参考文献}

\noindent 1.\quad Abraham, B. and J. Ledolter. \textit{Introduction to Regression Modeling}, Belmont, CA: Thompson Brooks/Cole, 2006.

\noindent 2.\quad Anderson, T. W. \textit{An Introduction to Multivariate Statistical Analysis} (3rd ed.). New York: John Wiley, 2003.
